% Algorithmic description of the RMS Double DQN pipeline
\documentclass{article}
\usepackage{algorithm}
\usepackage{algpseudocode}
\usepackage{amsmath}
\usepackage{geometry}
\geometry{margin=1in}

\begin{document}
\title{Reconfigurable Manufacturing System Double DQN Pipeline}
\author{Project Documentation}
\date{\today}
\maketitle

\begin{abstract}
This document presents an end-to-end algorithmic description, in pseudocode form, of the Reconfigurable Manufacturing System (RMS) scheduling pipeline that integrates real and synthetic data, reward shaping, and advanced Deep Q-Learning techniques. Each stage is detailed to facilitate reproducibility and scientific communication.
\end{abstract}

\begin{algorithm}
\caption{RMS Double DQN Training and Evaluation Pipeline}
\begin{algorithmic}[1]
\Require Configuration dictionary $\mathcal{C}$ specifying paths, hyperparameters, reward weights, exploration schedule, and visualization settings.
\Require Real CSV dataset $D_{\text{real}}$, synthetic generation budget $N_{\text{syn}}$, number of training episodes $E$, evaluation episodes $E_{\text{eval}}$.
\Ensure Trained policy network $Q_{\theta}$, diagnostics, and visual analytics.

\Statex
\Procedure{PrepareData}{$\mathcal{C}, D_{\text{real}}, N_{\text{syn}}$}
    \State Load $D_{\text{real}}$; remove inconsistent entries, impute missing values, and normalize numerical fields.
    \State Engineer RMS-specific features (critical ratio, slack, energy intensity, setup complexity, reconfiguration frequency).
    \State Fit parametric distributions to the cleaned dataset; sample $N_{\text{syn}}$ synthetic jobs respecting machine capabilities and deadline distributions.
    \State Return aggregated dataset $D = D_{\text{real}} \cup D_{\text{syn}}$ and normalization statistics.
\EndProcedure

\Statex
\Procedure{InitializeAgent}{$\mathcal{C}, D$}
    \State Instantiate RMS environment with state encoder using statistics from $D$.
    \State Build dueling Double DQN networks $(Q_{\theta}, Q_{\bar{\theta}})$ with architecture given by $\mathcal{C}$.
    \State Initialize prioritized experience replay buffer with capacity $\mathcal{C}[\text{buffer\_size}]$.
    \State Configure exploration schedules (\textit{e.g.}, $\varepsilon$-greedy decay, Boltzmann temperature) and reward shaping coefficients.
\EndProcedure

\Statex
\Procedure{Train}{$E$}
    \For{episode $= 1$ to $E$}
        \State Reset environment to sampled RMS scenario; receive initial state $s_0$.
        \For{each decision step $t$}
            \State Sample action $a_t$ from exploration policy using $Q_{\theta}(s_t)$ and valid-action mask.
            \State Execute $a_t$; observe reward $r_t$, next state $s_{t+1}$, feasibility flag $f_t$.
            \State Compose shaped reward $r_t^{\prime}$ combining tardiness, energy, setup, idle penalties, and throughput bonuses per $\mathcal{C}$.
            \State Store transition $(s_t, a_t, r_t^{\prime}, s_{t+1}, f_t)$ with priority $|\delta_t|$ in replay buffer.
            \If{buffer size $> \mathcal{C}[\text{batch\_size}]$}
                \State Sample mini-batch with prioritized importance weights.
                \State Compute Double DQN targets using $Q_{\theta}$ for action selection and $Q_{\bar{\theta}}$ for evaluation.
                \State Update $Q_{\theta}$ via gradient descent with clipped gradients and learning rate scheduler.
                \State Update replay priorities with new TD-errors.
            \EndIf
            \State Apply scheduled decay or warm restarts to exploration parameters.
        \EndFor
        \State Soft-update target network parameters $Q_{\bar{\theta}} \leftarrow \tau Q_{\theta} + (1-\tau) Q_{\bar{\theta}}$.
        \State Log episode metrics (returns, tardiness, energy consumption, buffer diagnostics).
    \EndFor
\EndProcedure

\Statex
\Procedure{Evaluate}{$E_{\text{eval}}$}
    \For{episode $= 1$ to $E_{\text{eval}}$}
        \State Run greedy policy derived from $Q_{\theta}$ on held-out RMS scenarios.
        \State Record performance metrics: throughput, lateness distribution, energy usage, reconfiguration count, machine utilization fairness.
    \EndFor
    \State Aggregate statistics and compare against baseline heuristics if available.
\EndProcedure

\Statex
\Procedure{VisualizeAndReport}{ }
    \State Generate 30 analytics plots summarizing training curves, exploration dynamics, Q-value stability, RMS KPIs, and scheduling timelines.
    \State Persist trained weights, configuration snapshot, and serialized diagnostics for reproducibility.
\EndProcedure

\Statex
\State \textbf{return} $Q_{\theta}$, evaluation report, and visualization assets.
\end{algorithmic}
\end{algorithm}

\end{document}
